% -------------------------------------------------------
%  English Abstract
% -------------------------------------------------------

\begin{latin}

\begin{center}
\textbf{Abstract}
\end{center}
\baselineskip=.8\baselineskip

The mismatch between conventional training data and real telephone and Voice over Internet Protocol (VoIP) speech is a major obstacle to robust Persian automatic speech recognition. Environmental noise, bandwidth limitations, codec compression, and packet loss can alter acoustic cues, while collecting and transcribing real Persian telephone calls is costly and raises privacy concerns. Pretrained speech models partly alleviate data scarcity, but large-scale pretraining does not eliminate the need for adaptation to the target language and channel. Prior work has investigated training with degraded data, speech enhancement, multi-view fusion, and the use of unlabeled audio as individual approaches to improving robustness. However, many of these approaches have not yet been studied extensively and systematically for Persian speech recognition, particularly under telephone and VoIP conditions.

This thesis aims to design and evaluate a practical, reproducible approach to robust Persian speech recognition without relying on the large-scale collection of labeled telephone calls. To this end, two routes are pursued. The first is data-driven and seeks to improve model robustness by simulating telephone and VoIP conditions in the training data and by increasing the amount of Persian speech available for training. The second is model-driven and examines whether an enhancement--fusion architecture can extract complementary information from noisy and enhanced views beyond what is already learned through multi-condition training. The principal innovation is to investigate these approaches within a common experimental framework, allowing the effects of training with telecommunications degradations, using a more complex architecture, and adding more data to be examined separately, while determining the extent to which improvements under synthetic conditions persist in real telephone speech.

In the data-driven route, a set of traceable telecommunications degradations is applied to Persian speech, and Whisper-small is trained under multiple conditions using a mixture of conventional and degraded speech. As a complementary approach, Iran Seda audiobooks are segmented, pseudo-labeled, filtered, and then added to the most complete training mixture to determine how increasing the amount of data, even without a direct match to the telephone channel, affects model performance. In the model-driven route, a dual-stream system is developed to process the original and enhanced views simultaneously and combine them through a learned mechanism. All principal routes are evaluated under a common protocol on several Persian datasets covering both general-domain speech and real telephone and VoIP calls.

The results show that multi-condition training substantially reduces overall recognition error without modifying the recognizer architecture, and that this improvement also transfers to real telecommunications speech. The model-driven route performs better than conventional training but does not surpass the multi-condition baseline; the behavior of its learned gate further shows that the model relies almost exclusively on the original view and does not obtain stable complementary information from the enhanced view. Adding pseudo-labeled speech, while maintaining approximately the same performance as before on real calls, improves the overall results and model performance on several general-domain datasets.

Based on this evidence, multi-condition Whisper-small is proposed as the selected model for the target application. Without adding architectural complexity or using transcribed telephone calls in training, it provides a clear improvement on real telecommunications speech. The central finding of this thesis is that, for this low-resource problem, matching the training data to the target channel is more consequential than directly adding architectural components or merely increasing the amount of training data.

\bigskip\noindent\textbf{Keywords}:
Persian automatic speech recognition, domain mismatch, telephone speech, Voice over IP, multi-condition training, telecommunications data augmentation, pseudo-labeling, dual-view fusion

\end{latin}
